\magicSection{Procedures \& Instructions}{sec:Proc}

Each procedure (including instructions) takes control variables and window references (Section~\ref{sec:Windows}) as parameters.

For static safety checking, the sub-procedure has
\begin{itemize}
  \item a required device scope, CPU or CUDA (Section~\ref{sec:SeqPar})
  \item a collective unit $\tau_u$ (e.g. warp); calls to the proc must appear at $\tau_u$-scope (Section~\ref{sec:SeqPar})
  \item memory types for each window reference (Section~\ref{sec:Tensors})
  \item typed strides for each dimension of each window reference; the real window references passed must have the expected strides (Section~\ref{sec:Windows}) and also swizzle function
\end{itemize}

Certain instructions may require window references to be passed in a special binary format (e.g. CUtensorMap, wgmma SMEM descriptor).

For what will initially be dynamic safety checking, the sub-procedure has
\begin{itemize}
  \item Value specification: each modified window has the expected output given as an extended polynomial function (Section~\ref{sec:PolynomialSpec}) of index coordinates and input buffers
  \item Pre-conditions and post-conditions of some sort on the synchronization environment (Section~\ref{sec:SyncCheck})
\end{itemize}

The sub-procedure call will have to have each window expression annotated in some way similar to assignment statements (Section~\ref{sec:AssignmentAnnotation}).

Unlike Exo-GPU, we should allow non-instruction procs that require CUDA scope.
This is because breaking giant kernels into sub-procedures may be very useful to cut down the complexity of proofs and provide a natural boundary of where to declare pre-conditions and post-conditions.
We will NOT translate these to CUDA C++ \lighttt{\_\_device\_\_} functions; instead, we should inline all such function calls inside the compiler so we don't have to trust the downstream compiler.
This is especially important for register windows, so we don't have to try to pass pointers/references to values that fundamentally are not stored in RAM.
